\section{Lecture 7: The Solid Affine Line}
\label{lec:7}

Clausen resumes the course and begins to do geometry. The object of study is
the \emph{solid affine line}: the polynomial ring $\mathbb{Z}[T]$, viewed
through the lens of solid modules, and the various subspaces of it (open and
closed unit disks) that the solid formalism cuts out. The lecture has two
halves. First, a relative version of solidity over $\mathbb{Z}[T]$ is set up,
its structure theory is proved by transporting Scholze's arguments through an
explicit formula for the relative solidification, and the resulting
categories are shown to fit into a \emph{recollement}
--- from which Clausen reads off a definite verdict on what is ``open'' and
what is ``closed'' in nonarchimedean geometry. Second, for a general solid
ring $R$ the theory is used to define the notions of \emph{power-bounded} and
\emph{topologically nilpotent} elements, to build the \emph{rational
localizations} $R\langle f_1,\dots,f_n/g\rangle$, and to compare the result
with Huber's theory --- of which it is a derived, non-quasiseparated
refinement that repairs the failure of Huber's structure presheaf to be a
sheaf.

Throughout, ``ring'' means commutative ring, and everything takes place inside
light condensed abelian groups. We freely use the theory of solid abelian
groups $\Solid\subset\CondAbl$ from \cref{lec:5,lec:6}.

\subsection{Recap}

Recall from \cref{lec:5} the abelian category
\begin{equation}
\label{eq:7:solid-recall}
\Solid\ \subset\ \CondAbl
\end{equation}
of solid abelian groups: a full subcategory, closed under limits, colimits,
extensions and much else, which one thinks of as an algebraically well-behaved
replacement for complete nonarchimedean topological abelian groups
(\cref{thm:6:recall-structure}). The inclusion has a colimit-preserving left
adjoint, \emph{solidification} $M\mapsto M^{\solid}$, and a symmetric monoidal
structure $\otimes^{\solid}$ (a completed tensor product) for which
solidification is symmetric monoidal (\cref{cor:5:solidification}). The
compact projective generator is $\prod_{\mathbb{N}}\mathbb{Z}$.

Solidity was defined (\cref{def:5:solid}) through the free light condensed
abelian group on a null sequence
\begin{equation}
\label{eq:7:P-recall}
P=\mathbb{Z}[\mathbb{N}\cup\{\infty\}]\big/\mathbb{Z}[\infty],
\end{equation}
which is internally projective (\cref{thm:3:three-properties}(3)): an object
$M$ is solid iff the endomorphism $\mathrm{id}-\mathrm{shift}$ of the shift on
$\mathbb{N}$ induces an isomorphism
\begin{equation}
\label{eq:7:solid-def-recall}
\iHom(P,M)\xrightarrow{\ 1-\sigma\ }\iHom(P,M),
\end{equation}
i.e.\ every null sequence in $M$ is uniquely summable.

\subsection{The ring \texorpdfstring{$P$}{P} and its solidification}

The new observation is that $P$ is not merely an abelian group but a
\emph{ring}, and that there is a ring map
\begin{equation}
\label{eq:7:ZT-to-P}
\mathbb{Z}[T]\longrightarrow P,\qquad T\longmapsto\sigma,
\end{equation}
under which the variable $T$ acts as the (injective) shift $\sigma$ that
enters \eqref{eq:7:solid-def-recall}.

\begin{remark}[The ring structure on $P$, via measures]
\label{rmk:7:P-ring}
The ring structure is cleanest when one attaches the measure construction to
$\mathbb{N}$ rather than to $\mathbb{N}\cup\{\infty\}$. For any countable set
$S$ one defines the measures on $S$ as the free module on any compactification
modulo the free module on the boundary; the result is independent of the
compactification, is functorial for proper maps, and is symmetric monoidal, in
that measures on $S_1$ tensored with measures on $S_2$ are measures on
$S_1\times S_2$ (seen using independence of the compactification). Since
addition $\mathbb{N}\times\mathbb{N}\to\mathbb{N}$ is a proper map with finite
fibres, it induces a multiplication on $P$, making $P$ a ring, and
\eqref{eq:7:ZT-to-P} is the induced $\mathbb{Z}[T]$-algebra structure.
\end{remark}

Solidifying $P$ is very clean. As an abelian group $P^{\solid}
=\prod_{\mathbb{N}}\mathbb{Z}$ (\cref{cor:5:P-solid}); taking the ring
structure into account, this is the one-variable power series ring:
\begin{equation}
\label{eq:7:P-solid-power-series}
P^{\solid}\ \cong\ \mathbb{Z}[\![T]\!].
\end{equation}
Thus, in the solid world, the universal carrier of a null sequence is the
power series ring $\mathbb{Z}[\![T]\!]$.

The following idempotence is the technical heart of the geometric picture.

\begin{lemma}[{$\mathbb{Z}[\![T]\!]$ is idempotent over $\mathbb{Z}[T]$}]
\label{lem:7:idempotent}
The multiplication map is an isomorphism
\begin{equation}
\label{eq:7:idempotent}
\mathbb{Z}[\![T]\!]\ \otimes^{\mathbb{L}\solid}_{\mathbb{Z}[T]}\
\mathbb{Z}[\![T]\!]\ \xrightarrow{\ \sim\ }\ \mathbb{Z}[\![T]\!],
\end{equation}
already at the level of the derived solid tensor product.
\end{lemma}

\begin{proof}
By the basic computation of the solid tensor product (\cref{cor:5:tensor-square}),
\begin{equation}
\label{eq:7:two-variable-power-series}
\mathbb{Z}[\![T_1]\!]\ \otimes^{\mathbb{L}\solid}_{\mathbb{Z}}\
\mathbb{Z}[\![T_2]\!]\ \cong\ \mathbb{Z}[\![T_1,T_2]\!].
\end{equation}
Forming \eqref{eq:7:idempotent} amounts to imposing $T_1=T_2$, i.e.\
quotienting by $T_1-T_2$, and $\mathbb{Z}[\![T_1,T_2]\!]/(T_1-T_2)
=\mathbb{Z}[\![T]\!]$. Alternatively, both sides are derived $T$-complete
(\cref{prop:6:tensor-p-complete}), so the isomorphism may be checked modulo
$T$, where it is a triviality.
\end{proof}

\begin{remark}[Derived solid tensor product]
\label{rmk:7:derived-tensor}
Here $\otimes^{\mathbb{L}\solid}$ is either the left derived functor of
$\otimes^{\solid}$, or equivalently the derived tensor product followed by
derived solidification. One need not know that there are enough flat objects to
form it, though in fact all the basic objects (products of copies of
$\mathbb{Z}[T]$) are flat (\cref{cor:6:flat}), so a resolution of one variable
suffices.
\end{remark}

\subsection{First interpretation: the open unit disk}

Because $\mathbb{Z}[\![T]\!]$ is idempotent over $\mathbb{Z}[T]$
(\cref{lem:7:idempotent}), inverting it inside $\mathbb{Z}[T]$-algebras does
not change it; such idempotent algebras correspond to \emph{subspaces} of
$\operatorname{Spec}\mathbb{Z}[T]$, i.e.\ of the affine line, rather than to
arbitrary spaces mapping to it. So $\mathbb{Z}[\![T]\!]$ should be thought of
as some subspace of $\mathbb{A}^1$.

The naive guess --- that $\mathbb{Z}[\![T]\!]$ is the \emph{formal completion
of $\mathbb{A}^1$ at the origin} --- is \emph{wrong}. The reason is that the
formal-neighbourhood interpretation is not stable under base change, whereas
$\mathbb{Z}[\![T]\!]^{\solid}$ is. One sees the difference after base change to
a nonarchimedean field.

\begin{example}[Base change to $\mathbb{Q}_p$: bounded functions on the open disk]
\label{ex:7:open-disk-Qp}
Base change $\mathbb{Z}[\![T]\!]$ to $\mathbb{Q}_p$. Since
$\mathbb{Q}_p=\mathbb{Z}_p[\tfrac1p]$ and the solid tensor product commutes
with colimits in each variable,
\begin{equation}
\label{eq:7:Qp-tensor-power-series}
\mathbb{Q}_p\ \otimes^{\mathbb{L}\solid}_{\mathbb{Z}}\ \mathbb{Z}[\![T]\!]
\ =\ \Bigl(\mathbb{Z}_p\ \otimes^{\mathbb{L}\solid}_{\mathbb{Z}}\
\mathbb{Z}[\![T]\!]\Bigr)\bigl[\tfrac1p\bigr].
\end{equation}
Now $\mathbb{Z}_p=\mathbb{Z}[\![U]\!]/(U-p)$ has a two-term resolution by
copies of the power series ring, so, using
\eqref{eq:7:two-variable-power-series} to compute
$\mathbb{Z}[\![U]\!]\otimes^{\mathbb{L}\solid}\mathbb{Z}[\![T]\!]
=\mathbb{Z}[\![U,T]\!]$ and pulling the (naive) limit through,
\begin{equation}
\label{eq:7:open-disk-answer}
\mathbb{Q}_p\ \otimes^{\mathbb{L}\solid}_{\mathbb{Z}}\ \mathbb{Z}[\![T]\!]
\ \cong\ \mathbb{Z}_p[\![T]\!]\bigl[\tfrac1p\bigr].
\end{equation}
This is \emph{not} the formal completion $\mathbb{Q}_p[\![T]\!]$; it is the
subring of $\mathbb{Q}_p[\![T]\!]$ of power series whose coefficients are
\emph{bounded} in $p$-adic norm. Its interpretation: it is the ring of bounded
rigid-analytic functions on the \emph{open} unit disk. For any nonarchimedean
extension of $\mathbb{Q}_p$, such a series converges exactly at points of
absolute value strictly less than $1$; all these functions converge on the
open unit disk and no further, but boundedness of the coefficients cuts out the
bounded functions rather than all holomorphic ones. Over a discrete ring such
as $\mathbb{Z}$ the open unit disk and the formal neighbourhood cannot be told
apart; the distinction emerges only after base change to a ring with more room
to move.
\end{example}

\subsection{The closed unit disk and \texorpdfstring{$\mathbb{Z}[T]$}{Z[T]}-solidity}

In rigid geometry the fundamental object is the \emph{closed} unit disk, not
the open one. To reach it, put the open unit disk ``at infinity'': working with
the variable $T^{-1}$, the ring $\mathbb{Z}[\![T^{-1}]\!]$ is the open unit
disk centred at $\infty$, and puncturing at $\infty$ to land in the affine line
gives the $\mathbb{Z}[T]$-algebra
\begin{equation}
\label{eq:7:complement}
\mathbb{Z}[\![T^{-1}]\!]\ \otimes^{\solid}_{\mathbb{Z}[T^{-1}]}\
\mathbb{Z}[T,T^{-1}]\ =\ \mathbb{Z}(\!(T^{-1})\!),
\end{equation}
the Laurent series in $T^{-1}$. This is the ring of the \emph{complement} of
the closed unit disk. To pass to the closed unit disk itself one must localize
away from this complement --- equivalently, \emph{kill} the object
$\mathbb{Z}(\!(T^{-1})\!)$.

As a $\mathbb{Z}[T]$-module, $\mathbb{Z}(\!(T^{-1})\!)$ has a simple two-term
resolution. Writing $U=T^{-1}$, one inverts $U$ inside $\mathbb{Z}[\![U]\!]$ by
adjoining a formal inverse and imposing the relation, giving
\begin{equation}
\label{eq:7:complement-resolution}
\Bigl[\ \mathbb{Z}[\![U]\!][T]\xrightarrow{\ UT-1\ }\mathbb{Z}[\![U]\!][T]\ \Bigr]
\ \xrightarrow{\ \sim\ }\ \mathbb{Z}(\!(T^{-1})\!).
\end{equation}
The two terms are the base change of $P$ to $\mathbb{Z}[T]$-modules in solid
abelian groups: $\mathbb{Z}[\![U]\!][T]=P^{\solid}\otimes_{\mathbb{Z}}
\mathbb{Z}[T]$ is the compact projective generator of $\mathbb{Z}[T]$-modules
in solid abelian groups. Thus $\mathbb{Z}(\!(T^{-1})\!)$ is resolved by
compact projective generators, and killing it should mean requiring the map
induced by $UT-1$ to become an isomorphism. By analogy with
\eqref{eq:7:solid-def-recall} we make the following definition.

\begin{definition}[{$\mathbb{Z}[T]$-solid modules}]
\label{def:7:ZT-solid}
Let $M$ be a module over $\mathbb{Z}[T]$ in $\CondAbl$, with $t\colon M\to M$
the action of $T$. Then $M$ is \emph{$\mathbb{Z}[T]$-solid} if
\begin{equation}
\label{eq:7:ZT-solid-def}
\iHom(P,M)\ \xrightarrow{\ \sigma t-1\ }\ \iHom(P,M)
\end{equation}
is an isomorphism, where $\sigma$ is the shift. Write $\Solid_{\mathbb{Z}[T]}$
for the full subcategory of $\mathbb{Z}[T]$-solid objects.
\end{definition}

By \eqref{eq:7:complement-resolution}, passing to $R\iHom$, the condition
\eqref{eq:7:ZT-solid-def} is equivalent to
\begin{equation}
\label{eq:7:ZT-solid-vanishing}
R\iHom_{\mathbb{Z}[T]}\bigl(\mathbb{Z}(\!(T^{-1})\!),\ M\bigr)\ =\ 0,
\end{equation}
i.e.\ to the demand that $M$ ``see nothing'' of the complement of the closed
unit disk.

\begin{theorem}[{Structure of the $\mathbb{Z}[T]$-solid theory}]
\label{thm:7:ZT-solid-structure}
Inside $\mathrm{Mod}_{\mathbb{Z}[T]}(\CondAbl)$, the category
$\Solid_{\mathbb{Z}[T]}$ satisfies:
\begin{enumerate}
\item it is an abelian subcategory, closed under limits, colimits and
extensions, and $\Solid_{\mathbb{Z}[T]}\subseteq\mathrm{Mod}_{\mathbb{Z}[T]}
(\Solid_{\mathbb{Z}})\subseteq\mathrm{Mod}_{\mathbb{Z}[T]}(\CondAbl)$;
\item if $M\in\mathrm{Mod}_{\mathbb{Z}[T]}(\CondAbl)$ is arbitrary and $N\in
\Solid_{\mathbb{Z}[T]}$, then $\iExt^{i}_{\mathbb{Z}[T]}(M,N)\in
\Solid_{\mathbb{Z}[T]}$ for all $i$;
\item the inclusion has a colimit-preserving left adjoint, the
\emph{$\mathbb{Z}[T]$-solidification} $(-)^{\Tsolid}$, and a unique symmetric
monoidal structure making it symmetric monoidal;
\item the derived analogue of \textup{(1)--(3)} holds, and the derived theory
is ``naive'': a complex is derived $\mathbb{Z}[T]$-solid iff each of its
homology groups is $\mathbb{Z}[T]$-solid, so that $D(\Solid_{\mathbb{Z}[T]})$
is the derived category of the abelian category $\Solid_{\mathbb{Z}[T]}$.
\end{enumerate}
\end{theorem}

The name is well chosen: the $\mathbb{Z}[T]$-solid theory over $\mathbb{Z}[T]$
is as robust as the solid theory over $\mathbb{Z}$.

\begin{remark}[{$\mathbb{Z}[T]$ and $\mathbb{Z}[\![T]\!]$ are $\mathbb{Z}[T]$-solid}]
\label{rmk:7:ZT-is-solid}
Both $\mathbb{Z}[T]$ and $\mathbb{Z}[\![T]\!]$ are $\mathbb{Z}[T]$-solid.
Indeed $\mathbb{Z}[T]$ is a retract of the compact projective generator (below),
and it is even solid as an abelian group, since every discrete abelian group is
solid (being generated under colimits by $\mathbb{Z}$). And $\mathbb{Z}[\![T]\!]$
is built from $\mathbb{Z}[T]$ by limits and colimits, hence is
$\mathbb{Z}[T]$-solid too.
\end{remark}

\begin{remark}[Nonarchimedean interpretation]
\label{rmk:7:nonarch-interpretation}
Just as $\Solid_{\mathbb{Z}}$ is an algebraic surrogate for complete
nonarchimedean topological abelian groups --- those with a basis of
neighbourhoods of $0$ consisting of subgroups --- the category
$\Solid_{\mathbb{Z}[T]}$ is a surrogate for complete nonarchimedean objects
admitting a basis of neighbourhoods of $0$ consisting of
$\mathbb{Z}[T]$-\emph{submodules}. The ring $\mathbb{Z}(\!(T^{-1})\!)$ is the
prototypical nonarchimedean $\mathbb{Z}[T]$-module for which no such basis
exists: its neighbourhood basis is stable under multiplication by $T^{-1}$
rather than by $T$. Killing this single object accounts for the whole
difference between the two notions.
\end{remark}

\subsection{An explicit formula for the
\texorpdfstring{$\mathbb{Z}[T]$}{Z[T]}-solidification}

The formal properties (1)--(3) of \cref{thm:7:ZT-solid-structure} are proved
exactly as Scholze proved their $\Solid_{\mathbb{Z}}$-analogues in
\cref{lec:5}: everything followed from $P$ being internally projective and from
asking that an endomorphism of $\iHom(P,-)$ become invertible, which is exactly
the shape of \eqref{eq:7:ZT-solid-def}. What must be redone is the
identification of the free objects; and here, thanks to the reading of the
theory as killing $\mathbb{Z}(\!(T^{-1})\!)$, matters are actually easier than
over $\mathbb{Z}$, because $\Solid_{\mathbb{Z}}$ is already available.

The key is an explicit formula for the (derived) $\mathbb{Z}[T]$-solidification.

\begin{proposition}[Formula for $(-)^{\LTsolid}$]
\label{prop:7:solidification-formula}
For $M\in D(\mathrm{Mod}_{\mathbb{Z}[T]}(\CondAbl))$,
\begin{equation}
\label{eq:7:solidification-formula}
M^{\LTsolid}\ =\ R\iHom_{\mathbb{Z}[T]}\Bigl(\
\mathbb{Z}(\!(T^{-1})\!)\big/\mathbb{Z}[T]\ [-1],\ M\ \Bigr),
\end{equation}
and the natural map $M\to M^{\LTsolid}$ is induced by
$R\iHom_{\mathbb{Z}[T]}(\mathbb{Z}[T],M)\simeq M$ under the boundary map of the
fibre sequence $\mathbb{Z}[T]\to\mathbb{Z}(\!(T^{-1})\!)\to
\mathbb{Z}(\!(T^{-1})\!)/\mathbb{Z}[T]$.
\end{proposition}

\begin{proof}[Sketch]
This is formal from the idempotence of $\mathbb{Z}(\!(T^{-1})\!)$ over
$\mathbb{Z}[T]$,
\begin{equation}
\label{eq:7:complement-idempotent}
\mathbb{Z}(\!(T^{-1})\!)\ \otimes^{\mathbb{L}\solid}_{\mathbb{Z}[T]}\
\mathbb{Z}(\!(T^{-1})\!)\ \xrightarrow{\ \sim\ }\ \mathbb{Z}(\!(T^{-1})\!),
\end{equation}
which follows from \eqref{eq:7:complement} and \cref{lem:7:idempotent} by
transporting to $\infty$. The object $\mathbb{Z}(\!(T^{-1})\!)/\mathbb{Z}[T]$
is the homotopy fibre of $\mathbb{Z}[T]\hookrightarrow\mathbb{Z}(\!(T^{-1})\!)$
shifted, and \eqref{eq:7:complement-idempotent} says $R\iHom_{\mathbb{Z}[T]}
(\mathbb{Z}(\!(T^{-1})\!)/\mathbb{Z}[T][-1],-)$ is an idempotent operation
whose local objects are exactly the $M$ with
$R\iHom_{\mathbb{Z}[T]}(\mathbb{Z}(\!(T^{-1})\!),M)=0$, i.e.\ the
$\mathbb{Z}[T]$-solid ones. This is precisely the situation of a Bousfield
localization, and \eqref{eq:7:solidification-formula} is the localization
functor.
\end{proof}

Something better happens when $M$ is base-changed from a solid
$\mathbb{Z}$-module. This is the computation that identifies the free objects.

\begin{proposition}[Pushforward from $\Solid_{\mathbb{Z}}$]
\label{prop:7:pushforward}
The functor
\begin{equation}
\label{eq:7:pushforward}
D(\Solid_{\mathbb{Z}})\longrightarrow D(\Solid_{\mathbb{Z}[T]}),\qquad
M\longmapsto\bigl(M\otimes_{\mathbb{Z}}\mathbb{Z}[T]\bigr)^{\LTsolid},
\end{equation}
is $t$-exact, preserves all limits and colimits, and sends $\mathbb{Z}$ to
$\mathbb{Z}[T]$. Concretely, there is a natural isomorphism
\begin{equation}
\label{eq:7:pushforward-formula}
\bigl(M\otimes_{\mathbb{Z}}\mathbb{Z}[T]\bigr)^{\LTsolid}\ \cong\
R\iHom_{\mathbb{Z}}\bigl(\,U\,\mathbb{Z}[\![U]\!]\,[-1],\ M\,\bigr),
\end{equation}
where the residual $\mathbb{Z}[T]$-module structure is induced by the
endomorphism of $U\,\mathbb{Z}[\![U]\!]$ sending $U\mapsto0$, $U^2\mapsto U$,
$U^3\mapsto U^2$, and so on.
\end{proposition}

\begin{proof}[Sketch]
Apply \eqref{eq:7:solidification-formula} to $M[T]:=M\otimes_{\mathbb{Z}}
\mathbb{Z}[T]$ and compute the $R\iHom_{\mathbb{Z}[T]}$ by disambiguating the
two occurrences of $T$ into $U$ and $T$ and then equalizing. Using the two-term
presentation of $\mathbb{Z}(\!(T^{-1})\!)$ from
\eqref{eq:7:complement-resolution}, one is led to
\begin{equation}
\label{eq:7:solidification-intermediate}
R\iHom_{\mathbb{Z}}\bigl(U\mathbb{Z}[\![U]\!][-1],\ M\bigr)[T]
\xrightarrow{\ T-\sigma\ }
R\iHom_{\mathbb{Z}}\bigl(U\mathbb{Z}[\![U]\!][-1],\ M[T]\bigr),
\end{equation}
after which one equalizes $T$ against the shift $\sigma$ induced by
multiplication by $T$. Now $U\mathbb{Z}[\![U]\!]$ is a compact projective object
of $\Solid_{\mathbb{Z}}$, so $R\iHom$ out of it commutes with the colimit
$M[T]=\bigoplus_n M\cdot T^n$; the resulting fibre (from ``$T$ minus a shift'')
together with the extra shift cancel, and one is left with the clean formula
\eqref{eq:7:pushforward-formula}. Since $U\mathbb{Z}[\![U]\!]$ is internally
projective in $\Solid_{\mathbb{Z}}$, the functor is $t$-exact and preserves
limits and colimits; and limits and colimits in $\Solid_{\mathbb{Z}[T]}$ are
computed on underlying objects, being a module category. Finally, plugging
$M=\mathbb{Z}$ into \eqref{eq:7:pushforward-formula}, $R\iHom_{\mathbb{Z}}$ out
of the product $U\mathbb{Z}[\![U]\!]=\prod_{\mathbb{N}}\mathbb{Z}$ into
$\mathbb{Z}$ turns into a direct sum, and a natural comparison exhibits the
result as $\mathbb{Z}[T]$ with its evident $\mathbb{Z}[T]$-module structure.
\end{proof}

\begin{corollary}[The compact projective generator]
\label{cor:7:generator}
Solidifying the sequence space produces the free module of the theory: for a
product of copies of $\mathbb{Z}$,
\begin{equation}
\label{eq:7:generator}
\Bigl(\bigl(\textstyle\prod_{\mathbb{N}}\mathbb{Z}\bigr)\otimes_{\mathbb{Z}}
\mathbb{Z}[T]\Bigr)^{\Tsolid}\ \cong\ \prod_{\mathbb{N}}\mathbb{Z}[T],
\end{equation}
a product of copies of $\mathbb{Z}[T]$, and this is the compact projective
generator of $\Solid_{\mathbb{Z}[T]}$.
\end{corollary}

\begin{proof}
The functor of \cref{prop:7:pushforward} commutes with products and sends
$\mathbb{Z}$ to $\mathbb{Z}[T]$; apply it to $\prod_{\mathbb{N}}\mathbb{Z}$.
\end{proof}

\begin{remark}[Where the derived solidification fails to be $t$-exact]
\label{rmk:7:not-t-exact}
The pushforward \eqref{eq:7:pushforward} is $t$-exact, but the
$\mathbb{Z}[T]$-solidification $(-)^{\LTsolid}$ on all of
$D(\mathrm{Mod}_{\mathbb{Z}[T]})$ is \emph{not}. By
\eqref{eq:7:solidification-formula} it is $R\iHom$ against an object with a
two-term compact resolution, so it is off from $t$-exactness by at most one:
connective objects stay connective, and coconnective objects go to at most a
one-fold shift of coconnective. It is thus very controlled, but genuinely not
exact. The pushforward \eqref{eq:7:pushforward}, by contrast, is $t$-exact even
though it factors through $(-)^{\LTsolid}$: it is the composite of the
$t$-exact base change $M\mapsto M\otimes_{\mathbb{Z}}\mathbb{Z}[T]$ with the
symmetric monoidal quotient
$\mathrm{Mod}_{\mathbb{Z}[T]}(\Solid_{\mathbb{Z}})\to\Solid_{\mathbb{Z}[T]}$,
and although this second, quotient step is not $t$-exact in general, on
base-changed objects the composite is.
\end{remark}

\begin{example}[Base change to $\mathbb{Q}_p$: the Tate algebra]
\label{ex:7:closed-disk-Qp}
Take $\mathbb{Q}_p[T]$, the functions on $\mathbb{A}^1_{\mathbb{Q}_p}$, and
restrict to the closed unit disk by applying $(-)^{\Tsolid}$, an instance of
\cref{prop:7:pushforward}. Pulling $\tfrac1p$ and the limit through,
\begin{equation}
\label{eq:7:closed-disk-computation}
\mathbb{Q}_p[T]^{\Tsolid}\ =\ \varprojlim_{n}\Bigl((\mathbb{Z}/p^n)[T]
\Bigr)^{\Tsolid}\bigl[\tfrac1p\bigr]\ =\ \varprojlim_{n}(\mathbb{Z}/p^n)[T]
\bigl[\tfrac1p\bigr]\ =\ \mathbb{Z}[T]^{\wedge}_{p}\bigl[\tfrac1p\bigr]
\ =\ \mathbb{Q}_p\langle T\rangle,
\end{equation}
using that $(\mathbb{Z}/p^n)[T]$ is base-changed from a discrete ring, so its
$\mathbb{Z}[T]$-solidification is itself. This is the Tate algebra: power series
in $T$ with coefficients in $\mathbb{Q}_p$ tending to $0$. The smaller ring of
functions converges on a larger region --- one is now allowed to sit on the
boundary $|T|=1$ --- exactly as the closed unit disk should.
\end{example}

\subsection{Open versus closed: a recollement}

We can now settle a question of terminology that nonarchimedean geometry
leaves ambiguous. In rigid geometry the closed unit disk is treated as an open
(it is affinoid, quasicompact), while the open unit disk is a non-quasicompact
increasing union of closed disks; from the Berkovich viewpoint the roles can
seem to switch again. Clausen's claim is that the categories attached to these
objects give a definite answer.

Consider the three symmetric monoidal derived categories: for the affine line
$\mathbb{A}^1$ the modules $D(\mathrm{Mod}_{\mathbb{Z}[T]}(\Solid_{\mathbb{Z}}))$;
for the closed unit disk the quotient $D(\Solid_{\mathbb{Z}[T]})$; and for the
complement of the closed unit disk (the open unit disk moved to $\infty$) the
modules $D(\mathrm{Mod}_{\mathbb{Z}(\!(T^{-1})\!)}(\Solid_{\mathbb{Z}}))$. They
sit in a diagram
% board 22 @ 00:44:46--00:52:47 — /home/deven/documents/coding/vms/gpu/home/notetaker/output/dustin-clausen-724-analytic-stacks/boards/board-22.jpg
\begin{equation}
\label{eq:7:recollement}
\begin{tikzcd}[column sep=large]
D\bigl(\mathrm{Mod}_{\mathbb{Z}(\!(T^{-1})\!)}(\Solid_{\mathbb{Z}})\bigr)
\arrow[r, "i_{*}"]
& D\bigl(\mathrm{Mod}_{\mathbb{Z}[T]}(\Solid_{\mathbb{Z}})\bigr)
\arrow[r, "j^{*}=(-)^{\LTsolid}"]
& D\bigl(\Solid_{\mathbb{Z}[T]}\bigr)
\end{tikzcd}
\end{equation}
which is formally a \emph{recollement} (localization sequence), realized by the
idempotent algebra $\mathbb{Z}(\!(T^{-1})\!)$ inside the symmetric monoidal
category of $\mathbb{Z}[T]$-modules. The two sides carry the standard
adjunctions.

\begin{itemize}
\item \emph{The closed unit disk, as an open $j$.} The symmetric monoidal
functor $j^{*}=(-)^{\LTsolid}$ has a right adjoint $j_{*}$, namely the
inclusion $D(\Solid_{\mathbb{Z}[T]})\hookrightarrow
D(\mathrm{Mod}_{\mathbb{Z}[T]})$; it is fully faithful because the idempotence
\eqref{eq:7:complement-idempotent} forces at most one
$\mathbb{Z}(\!(T^{-1})\!)$-module structure on a $\mathbb{Z}[T]$-module, so
$D(\Solid_{\mathbb{Z}[T]})$ is a symmetric monoidal quotient. It also has a
left adjoint $j_{!}$, given by
\begin{equation}
\label{eq:7:j-lower-shriek}
j_{!}\colon M\longmapsto\mathrm{fib}\Bigl(M\to M\otimes_{\mathbb{Z}[T]}
\mathbb{Z}(\!(T^{-1})\!)\Bigr),
\end{equation}
i.e.\ ``base change to $\infty$ and take the fibre''. This $j_{!}$ is better
behaved than $j_{*}$: it satisfies the projection formula over the two
symmetric monoidal categories, and is $j^{*}$-linear. As at infinity this is
extension by zero --- one keeps the sections and kills those living near the
boundary.
\item \emph{The complement, as a closed $i$.} The inclusion
$i_{*}\colon D(\mathrm{Mod}_{\mathbb{Z}(\!(T^{-1})\!)})\hookrightarrow
D(\mathrm{Mod}_{\mathbb{Z}[T]})$ is fully faithful; its left adjoint is the
symmetric monoidal base change $i^{*}=(-)\otimes_{\mathbb{Z}[T]}
\mathbb{Z}(\!(T^{-1})\!)$, and its right adjoint $i^{!}$ is a certain
$R\iHom$, which is less well-behaved. Here $i_{*}$ satisfies the projection
formula over $i^{*}$.
\end{itemize}

This is formally identical to the six-functor package attached to a
topological space $X$ with a closed subset $Z$ and open complement $U$:
\begin{equation}
\label{eq:7:topological-analogy}
\begin{tikzcd}[column sep=large]
D(\mathrm{Shv}(Z)) \arrow[r, "i_{*}"] & D(\mathrm{Shv}(X))
\arrow[r, "j^{*}"] & D(\mathrm{Shv}(U)),
\end{tikzcd}
\end{equation}
with $i^{*}\dashv i_{*}\dashv i^{!}$ and $j_{!}\dashv j^{*}\dashv j_{*}$, the
well-behaved (projection-formula) functors being $i_{*}$ and $j_{!}$. Both
\eqref{eq:7:recollement} and \eqref{eq:7:topological-analogy} are instances of
the same abstract input: a symmetric monoidal category together with an
idempotent algebra --- here $\mathbb{Z}(\!(T^{-1})\!)$, there the pushforward
of the constant sheaf from $Z$.

\begin{remark}[The verdict: closed disks are open]
\label{rmk:7:verdict}
Comparing \eqref{eq:7:recollement} with \eqref{eq:7:topological-analogy}, the
closed unit disk $D(\Solid_{\mathbb{Z}[T]})$ occupies the slot of the
\emph{open} $U$, and the complementary open unit disk occupies the slot of the
\emph{closed} $Z$. So from this perspective the closed unit disk should be
thought of as open, and the open unit disk as closed. The same reasoning
applies in usual algebraic geometry: a quasicompact (e.g.\ distinguished
affine) Zariski open, being the locus where a function is inverted --- an
idempotent algebra --- should be thought of as \emph{closed}. Clausen offered
this as a genuine reorientation, to be justified later, and noted that these
recollements are exactly the sort of structure guiding the definition of
analytic stack: one takes seriously the principle that the derived categories
and the functors between them dictate what the geometry is, and it is in this
condensed/solid framework that one obtains six-functor formalisms on derived
categories of quasi-coherent sheaves in large generality, with proofs of e.g.\
Serre duality.
\end{remark}

\subsection{Vista: solid rings and the valuative spectrum}

The plan for the rest of the theory is to run the above over a general
\emph{solid ring} --- a commutative algebra object
\begin{equation}
\label{eq:7:solid-ring}
R\in\CAlg\bigl(\Solid_{\mathbb{Z}},\otimes^{\solid}\bigr)
\end{equation}
--- and to organize the resulting subspaces (closed disks, open disks, and
their many analogues coming from arbitrary functions) in a systematic way. The
mechanism is to base-change the pictures found over $\mathbb{Z}[T]$ along all
possible maps $\mathbb{Z}[T]\to R$, i.e.\ along all functions in $R$.

\begin{remark}[Localization along the valuative spectrum]
\label{rmk:7:valuative-spectrum}
The outcome is that the derived category
$D(\mathrm{Mod}_{R}(\Solid_{\mathbb{Z}}))$ --- even as a symmetric monoidal
category --- localizes along the \emph{valuative spectrum}
$\Spv(R(\ast))$ of the underlying discrete ring $R(\ast)$: the space of all
valuations of $R(\ast)$ (with no continuity or integrality restriction --- an
arbitrary valuation, given by a residue field and a valuation on it). This is a
spectral space, with a basis of quasicompact opens; the particularly nice
basis, analogous to distinguished affine opens, consists of the
\emph{rational opens}. A rational open is cut out by finitely many functions
$f_1,\dots,f_n$ and a further function $g$,
\begin{equation}
\label{eq:7:rational-open}
X\Bigl(\frac{f_1,\dots,f_n}{g}\Bigr)\ =\ \bigl\{\,|f_i|\le|g|\ \forall i,\
\ g\neq0\,\bigr\},
\end{equation}
i.e.\ one inverts $g$ (lands in the Zariski nonvanishing locus of $g$) and
then shrinks to where each $|f_i|\le|g|$. One attaches to such an open a
version of the solid theory: the $M\in D(\mathrm{Mod}_{R}(\Solid_{\mathbb{Z}}))$
such that multiplication by $g$ is an isomorphism on $M$, and such that for
each $i$
\begin{equation}
\label{eq:7:rational-open-condition}
\iHom(P,M)\xrightarrow{\ (f_i/g)\,\sigma-1\ }\iHom(P,M)
\end{equation}
is an isomorphism; in words, thinking of $g,f_i$ as maps
$\operatorname{Spec}R\to\mathbb{A}^1$, one requires $g$ to land in the Zariski
nonvanishing locus and each $f_i/g$ to land in the closed unit disk. This
exhibits a sheaf of symmetric monoidal categories on $\Spv(R(\ast))$, and in
particular a structure sheaf. The goal of the rest of the lecture is to make
this structure sheaf explicit and to begin comparing it with Huber's theory;
the full construction, by easy formal means, needs more language and is
deferred.
\end{remark}

\subsection{Power-bounded and topologically nilpotent elements}

Fix $R\in\CAlg(\Solid_{\mathbb{Z}})$ and an element $f\in R(\ast)$ of the
underlying discrete ring. It determines a ring map $\mathbb{Z}[T]\to R$,
$T\mapsto f$, i.e.\ the sequence $1,f,f^2,f^3,\dots$.

\begin{definition}[Topologically nilpotent, power-bounded]
\label{def:7:tn-pb}
With notation as above:
\begin{itemize}
\item $f$ is \emph{topologically nilpotent} if the ring map
$\mathbb{Z}[T]\to R$ factors (as a ring map) through $\mathbb{Z}[\![T]\!]
=P^{\solid}$;
\item $f$ is \emph{power-bounded} if $R$, viewed as a $\mathbb{Z}[T]$-module
via $T\mapsto f$, is $\mathbb{Z}[T]$-solid, i.e.\
$\iHom(P,R)\xrightarrow{f\sigma-1}\iHom(P,R)$ is an isomorphism.
\end{itemize}
Write $R^{\circ}\subseteq R(\ast)$ for the set of power-bounded elements and
$R^{\circ\circ}\subseteq R(\ast)$ for the set of topologically nilpotent
elements (Huber's notation).
\end{definition}

\begin{remark}[Comparison with the classical definitions]
\label{rmk:7:classical-tn}
For topological nilpotence the agreement with the classical notion is
immediate: since $T\mapsto f$ is a ring map, $T^n\mapsto f^n$, and factoring
through $\mathbb{Z}[\![T]\!]$ says exactly that the sequence $1,f,f^2,\dots$
extends to a null sequence --- topological nilpotence. The genuinely new
feature in the solid setting is that such a factorization, \emph{if it exists,
is unique}, by the idempotence of $\mathbb{Z}[\![T]\!]$: even though solid
abelian groups can be very non-Hausdorff, the limit built into solidity is
imposed to exist uniquely.
\end{remark}

\begin{remark}[A subtlety: factorization as rings]
\label{rmk:7:factorization-rings}
The factorization defining topological nilpotence must be taken in the category
of \emph{rings}. A factorization merely as condensed abelian groups, or even as
$\mathbb{Z}[T]$-modules, is weaker: one may have a null sequence
$1,f,f^2,\dots$ (a factorization of abelian groups) without $f$ being
topologically nilpotent in the ring sense. If, however, the target is
quasiseparated --- the condensed analogue of Hausdorff --- then by density the
factorization is automatically a ring map when it exists, and the two notions
agree.\todo{Clausen and an audience member left slightly open whether ``null
sequence in the module sense'' always coincides with topological nilpotence in
general (non-quasiseparated) targets; the transcript (00:29--00:31,
01:07--01:10) settles the quasiseparated case and flags the general one as
subtle.}
\end{remark}

\begin{lemma}[Huber-type structure of $R^{\circ}$ and $R^{\circ\circ}$]
\label{lem:7:Rzero}
$R^{\circ}\subseteq R(\ast)$ is an integrally closed subring, and
$R^{\circ\circ}\subseteq R^{\circ}$ is a radical ideal.
\end{lemma}

\begin{proof}[Proof that $R^{\circ}$ is a subring]
That $1\in R^{\circ}$ is the definition of solidity (put $f=1$). For closure
under the ring operations, suppose $f,g\in R^{\circ}$; we show $F(f,g)\in
R^{\circ}$ for any $F\in\mathbb{Z}[X,Y]$. Consider the ring map
$\mathbb{Z}[X,Y]\to R$, $X\mapsto f$, $Y\mapsto g$; the hypothesis is that $R$
is $\mathbb{Z}[X]$-solid and $\mathbb{Z}[Y]$-solid, and we want it to be
$\mathbb{Z}[T]$-solid along $T\mapsto F(f,g)$.
% board 31--32 @ 01:11:36--01:15:22
\[
\begin{tikzcd}[row sep=small, column sep=large]
& \mathbb{Z}[T] \arrow[d, "T\mapsto F"] \\
\mathbb{Z}[X,Y] \arrow[r, "{X\mapsto f,\ Y\mapsto g}"'] & R
\end{tikzcd}
\]
Resolve $R$ (which is solid as an abelian group) by direct sums of the compact
projective generators $\prod_{\mathbb{N}}\mathbb{Z}$, base-changed to
$\mathbb{Z}[X,Y]$-modules. Solidifying with respect to $X$ turns each term into
$\bigoplus\bigl(\prod_{\mathbb{N}}\mathbb{Z}[\![X]\!]\bigr)[Y]$ (by the
properties of $\mathbb{Z}[X]$-solidification, \cref{prop:7:pushforward}) without
changing $R$; solidifying then with respect to $Y$ gives
$\bigoplus\prod_{\mathbb{N}}\mathbb{Z}[\![X,Y]\!]$. Each such term is
$\mathbb{Z}[T]$-solid --- a colimit of limits of things solid over $\mathbb{Z}$,
discrete in $T$, hence solid along any $T\mapsto F(f,g)$ --- and
$\Solid_{\mathbb{Z}[T]}$ is closed under colimits, so $R$ is
$\mathbb{Z}[T]$-solid. Thus $F(f,g)\in R^{\circ}$.
\end{proof}

\begin{proof}[Proof of integral closedness]
Suppose $f\in R(\ast)$ satisfies a monic equation
\begin{equation}
\label{eq:7:integral-equation}
f^{n}+c_{n-1}f^{n-1}+\cdots+c_{0}=0,\qquad c_i\in R^{\circ}.
\end{equation}
Form the universal situation: over $\mathbb{Z}[X_0,\dots,X_{n-1}]$ adjoin a
variable $T$ and impose $T^{n}+X_{n-1}T^{n-1}+\cdots+X_0=0$, with the map to
$R$ sending $X_i\mapsto c_i$ and $T\mapsto f$. By hypothesis $R$ is solid over
each $\mathbb{Z}[X_i]$; running the same resolution as above (solidify in each
$X_i$), one finds $R$ resolved by direct sums of products of copies of the
ring $\mathbb{Z}[X_0,\dots,X_{n-1}][T]/(T^n+\cdots+X_0)$, which is finite over
$\mathbb{Z}[X_0,\dots,X_{n-1}]$; a finite (indeed Noetherian, so finitely
presented) module can be pulled inside the products, and the resulting terms are
solid over $\mathbb{Z}[T]$. Hence $R$ is $\mathbb{Z}[T]$-solid along
$T\mapsto f$, i.e.\ $f\in R^{\circ}$. The proofs that $R^{\circ\circ}$ is an
ideal and is radical are entirely analogous and are left as exercises in the
same style.\footnote{Board~33 writes the coefficients as $c_i\in
R^{\circ\circ}$, but this is a slip: the statement being proved is integral
closedness of $R^{\circ}$, which by definition concerns monic relations with
$c_i\in R^{\circ}$, and the proof uses precisely that each $c_i$ is
power-bounded (so that $R$ is $\mathbb{Z}[X_i]$-solid via $X_i\mapsto c_i$).
Clausen indeed said ``all $c_i$ are power-bounded'' (01:16:55). Transcribed here
as $R^{\circ}$ accordingly.}
\end{proof}

\subsection{Rational localizations and the structure sheaf}

We can now build the rings that will be the values of the structure sheaf on
the rational opens \eqref{eq:7:rational-open}.

\begin{theorem}[Rational localization]
\label{thm:7:rational-localization}
Let $R\in\CAlg(\Solid_{\mathbb{Z}})$ and $g,f_1,\dots,f_n\in R(\ast)$. Then
there exists an initial solid ring
\begin{equation}
\label{eq:7:rational-localization}
R\Bigl\langle\frac{f_1,\dots,f_n}{g}\Bigr\rangle^{\solid},\qquad
R\longrightarrow R\Bigl\langle\frac{f_1,\dots,f_n}{g}\Bigr\rangle^{\solid},
\end{equation}
such that in the target
\begin{enumerate}
\item $g$ is invertible, and
\item $f_i/g$ is power-bounded for all $i$.
\end{enumerate}
Explicitly, it is constructed as
\begin{equation}
\label{eq:7:rational-construction}
R\Bigl\langle\frac{f_1,\dots,f_n}{g}\Bigr\rangle^{\solid}\ =\
\Bigl(R[X_1,\dots,X_n]^{X_1^{\blacksquare},\dots,X_n^{\blacksquare}}\big/
(gX_1-f_1,\ \dots,\ gX_n-f_n)\Bigr)\bigl[\tfrac1g\bigr],
\end{equation}
where one first freely adjoins the solid variables $X_i$ (i.e.\ solidifies with
respect to each), then imposes $gX_i=f_i$ (making $X_i=f_i/g$ power-bounded),
then inverts $g$.
\end{theorem}

\begin{proof}
The construction \eqref{eq:7:rational-construction} manifestly has the correct
universal property. Adjoining a solid variable $X_i$ and solidifying is what
forces $X_i$ power-bounded (this ``does something'' when $R$ is nondiscrete:
over $\mathbb{Q}_p$ with one variable it produces the Tate algebra, cf.\
\cref{ex:7:closed-disk-Qp}); imposing $gX_i-f_i$ then sets $X_i=f_i/g$; and
inverting $g$ makes $g$ a unit. The only point to check is that after these
operations the result is still solid, which holds because they are all
colimits.
\end{proof}

\begin{remark}[Two warnings: derived and non-quasiseparated]
\label{rmk:7:warnings}
Two cautions about \eqref{eq:7:rational-construction}.
\begin{enumerate}
\item It is \emph{not} exactly the value of the structure sheaf on
$X(f_1,\dots,f_n/g)$; it is $\pi_0$ of that value. The six-functor formalisms
and the localization of \cref{rmk:7:valuative-spectrum} live at the
\emph{derived} level, so what is produced a priori is a \emph{derived} sheaf.
In almost all practical cases the higher homotopy vanishes, $\pi_i=0$ for
$i>0$, so this causes no trouble; but it must be kept in mind.
\item Even for very nice $R$ --- a complete Huber ring --- the quotient in
\eqref{eq:7:rational-construction} need not be quasiseparated. The solidified
polynomial ring
$R[X_1,\dots,X_n]^{X_1^{\blacksquare},\dots,X_n^{\blacksquare}}$ is the usual
Tate algebra in several variables (this follows from the arguments above for
any complete Huber ring $R$), but the ideal generated by the $gX_i-f_i$ need
not be closed, so the quotient can be non-Hausdorff. Again, in essentially all
practical cases it is quasiseparated.
\end{enumerate}
\end{remark}

\begin{remark}[Relation to Huber's theory]
\label{rmk:7:huber}
Let $R$ be a complete Huber ring and suppose the ideal $(f_1,\dots,f_n)\subseteq
R$ is open --- the condition that describes rational opens in the space of
\emph{continuous} valuations, as opposed to all valuations. Then Huber
\cite{Huber_1994} defines a ring $R\langle f_1,\dots,f_n/g\rangle^{\mathrm{Huber}}$.
It is recovered from the solid construction as its \emph{quasiseparation}:
\begin{equation}
\label{eq:7:huber-comparison}
R\Bigl\langle\frac{f_1,\dots,f_n}{g}\Bigr\rangle^{\mathrm{Huber}}\ =\
\text{quasiseparation of}\ \pi_0\, R\Bigl\langle\frac{f_1,\dots,f_n}{g}
\Bigr\rangle^{\solid}.
\end{equation}
So Huber's ring is obtained functorially from the (more generally defined,
possibly derived and non-quasiseparated) solid object, but the two need not be
equal in general --- though in all practical cases they are. This repairs a
known defect of Huber's theory: in general his structure presheaf is only a
presheaf, not a sheaf (though it is a sheaf in all practical cases). If one
drops the insistence on quasiseparatedness and on being non-derived, the solid
theory yields a genuine structure sheaf --- indeed a derived category of
quasi-coherent sheaves that localizes --- and moreover allows one to define
these objects even without the openness hypothesis on $(f_1,\dots,f_n)$.
\end{remark}

\subsection{Closing remarks}

\begin{remark}[Non-continuous valuations and Colmez's ad hoc spaces]
\label{rmk:7:colmez}
Asked whether non-continuous valuations are ever used, Clausen pointed to
Colmez, who considered spaces (which he called ``ad hoc spaces'') designed to
include objects such as the open unit disk and $\mathbb{Z}_p$ as honest affine
objects, in the context of his work on the mod-$p$ Langlands correspondence for
$\mathrm{GL}_2$. The open unit disk is the prototypical topological ring that is
not a Huber ring, so in the usual rigid theory one is forced to present it as a
non-quasicompact union of closed disks; Colmez advocated sometimes treating it
directly as affine, which the present theory accommodates easily. Clausen noted
this was done in an example rather than as a developed theory.\todo{``Colmez'',
the term ``ad hoc spaces'', and the attribution to his mod-$p$
$\mathrm{GL}_2$ Langlands work are from the closing Q\&A (01:30:28--01:32:10),
stated tentatively; a precise reference was not given and is not supplied here.}
\end{remark}

\begin{remark}[Induced analytic ring structures; recovering Huber pairs]
\label{rmk:7:analytic-ring}
Every algebra $R$ in solid $\mathbb{Z}$-modules gives what will be called an
\emph{induced} analytic ring structure: its module category sits inside
condensed $R$-modules and forms an analytic ring in the sense of
\cref{def:1:analytic-ring}. For $R=\mathbb{Z}[T]$ the $\mathbb{Z}[T]$-solid
theory is a further, more complete such structure; both exist and one wants
both. This extra flexibility can be extended to any discrete ring, though for a
general solid $\mathbb{Z}$-algebra Clausen knew of no completely canonical
choice --- one option is to force all power-bounded elements to be solid, giving
the maximally complete structure obtainable with the present tools, though
further completions may exist.

Feeding a \emph{pair} $(R,R^{+})$ into the analytic-ring formalism of
\cref{lec:1}, rather than a bare ring, recovers an analogue of Huber's theory
for Huber pairs: one imposes solidity of the $f_i/g$ only for the chosen ring
of integral elements $R^{+}$, not for all functions. This fits Huber's theory
remarkably well, and topologically nilpotent elements are automatically solid
in any case.
\end{remark}

\begin{remark}[Solidifications commute]
\label{rmk:7:solidifications-commute}
A structural bonus of the present setup: whereas in the earlier theory forming
the tensor product of analytic rings required solidifying in one direction and
then the other (a colimit of iterated completions, cf.\
\cref{rmk:1:pushout-completion}), here the various solidifications all commute
with one another, since each is given by an $R\iHom$ out of a fixed object and
any two such $R\iHom$'s commute.
\end{remark}
